We define geometric intersections before ideal moments. On an actual primary blowup \(Y=\operatorname{Bl}_I\operatorname{Spec}R\), G331 supplies actual effective Cartier \(E_I\), ample \(\mathcal O_Y(-E_I)\), a reduced model whose components dominate the minimal components, and the exceptional Cartier order unit. Product-primary blowups give common refinements with actual invertible image ideals. Put \(p=n-1\).

For an effective exceptional integral Cartier \(E\), a power of \(\mathfrak m\) kills \(\mathcal O_E\): its coherent structure sheaf has closed-fibre support, so this follows on a finite affine cover. Its proper cohomology has finite \(R\)-length. Define
\[
\chi_R(E,\mathcal F)=\sum_j(-1)^j\ell_R H^j(E,\mathcal F).
\]
For line bundles \(L_1,\ldots,L_r\) on \(Y\), \(\chi_R(E,\mathcal O_E\otimes\bigotimes L_i^{t_i})\) is a polynomial of total degree at most \(p\). Indeed filter \(\mathcal O_E\) by \(\mathfrak m\)-powers. Its finitely many graded coherent quotients are on the proper \(k\)-scheme \(E_0\), of support dimension at most \(p\). Tensoring is exact; Euler additivity and complete Stacks 33.45.1 give the polynomial, with \(R\)-length equal to \(k\)-dimension on these quotients. This works in mixed characteristic without a coefficient field, and retains every nilpotent layer of \(E\).

For integral exceptional \(F_1,\ldots,F_p\) define \(Q_Y(F_1,\ldots,F_p,E)\) as the coefficient of \(t_1\cdots t_p\) in that polynomial with \(L_i=\mathcal O_Y(F_i)\). If \(p=0\), it is \(\ell_R H^0(E,\mathcal O_E)\): the supported proper zero-dimensional scheme is finite over its Artinian base and has no higher cohomology. Complete 33.45.2 applied to the filtration yields
\[
Q_Y(F_1,\ldots,F_p,E)
=\sum_{\dim Z=p}\ell_{\mathcal O_{Y,\eta_Z}}(\mathcal O_{E,\eta_Z})
 (F_1|Z\cdots F_p|Z\cdot Z). \tag{C}
\]
Here \(Z\) is integral reduced and proper over \(k\), and the parentheses are the numerical Euler-characteristic intersections of 33.45.3. At each top generic point, the filtration lengths add to the EXACT displayed local length: the residue field is unchanged by the \(\mathfrak m\)-quotient. Lower-dimensional filtration supports contribute zero. Thus all generic thickness weights are present. A fixed extra twist changes no degree-\(p\) coefficient.

Polynomial substitution proves additivity and symmetry of the first slots (33.45.5). The sequence
\[
0\longrightarrow\mathcal O_{E_1}(-E_2)
\longrightarrow\mathcal O_{E_1+E_2}
\longrightarrow\mathcal O_{E_2}\longrightarrow0
\]
and fixed-twist invariance prove additivity in the effective last slot. The Cartier order unit expresses every integral exceptional \(F\) as a difference of effective ones. This extends the last slot independently of decomposition, since equality of two differences is equality of their effective sums. Extend all slots to real Cartier divisors by multilinearity.

Full symmetry uses an ACTUAL supported Koszul complex. For effective exceptional \(E,F\), take \(K_E=[\mathcal O(-E)\to\mathcal O]\) and \(K_F=[\mathcal O(-F)\to\mathcal O]\), using their Cartier sections. Each is a locally free resolution of its quotient. Their tensor totals are identical up to the usual sign interchange. Tensor also with any line bundle \(L\). Compute its Euler characteristic on \(E\) and on \(F\); all terms and hypercohomology have finite-length cohomology, since support is proper over an Artinian base. The spectral-sequence Euler identity gives
\[
\chi_R(E,L)-\chi_R(E,L(-F))
=\chi_R(F,L)-\chi_R(F,L(-E)). \tag{K}
\]
Neither divisor must cut regularly on the other. Put
\(L=\mathcal O(F)^t\mathcal O(E)^s\bigotimes_{j=1}^{n-2}\mathcal O(G_j)^{r_j}\).
Set \(t=s=0\) and take the \(r_1\cdots r_{n-2}\) coefficient. The finite difference on the left extracts the \(t r_1\cdots r_{n-2}\) coefficient: total degree at most \(n-1\) excludes higher \(t\)-powers. The right extracts its \(s\) counterpart. Thus
\[
Q_Y(F,G_1,\ldots,G_{n-2},E)=Q_Y(E,G_1,\ldots,G_{n-2},F).
\]
Effective differences and real linearity give every transposition. For \(n=1\) there is no exchange. This pairing is supported locally, not a degree on global \(\mathrm{CH}_0(Y)\): for a DVR, the parameter divisor has \(Q=1\), although its global Chow class is zero.

For projection use \(\delta(P)=\dim R/P\). Catenarity and locality make this a dimension function; on finite-type models use its usual transcendence-degree extension. Excellence supplies the universally catenary hypotheses of complete Stacks 42.7.1. Top components have \(\delta\)-dimension \(n\), and closed-fibre dimension is ordinary \(k\)-dimension.

Let \(f:Y'\to Y\) be a primary refinement, finite normalization, or their normal-component lift. It is proper birational on each top component. Pullback of exceptional Cartier equations is regular, since these reduced components have generic points outside their support. At a dimension-\(p\) generic point of an exceptional divisor, only top components can occur: a smaller component dominates its own minimal component away from \(\mathfrak m\), and cannot contain a dimension-\(p\) closed-fibre subvariety. The local ring there is reduced one-dimensional. Embed such a local ring \(B\) into \(\bigoplus_{P\in\operatorname{Min}B}B/P\); the cokernel has finite length. Multiplication by a regular Cartier equation \(g\) has equal kernel and cokernel lengths on that finite-length module, giving
\[
\ell(B/gB)=\sum_P\ell((B/P)/g(B/P)).
\]
Thus the top divisor cycle is the sum of divisor cycles on the integral top components, whose generic reduced weights are one. Complete Stacks 42.26.3 applies to each component and gives the ACTUAL cycle equality
\[
f_*[f^*E]_p=[E]_p. \tag{P}
\]
Its function-field degree is one; residue-field degrees in cycle pushforward remain. Numerical projection, complete 33.45.7, on the proper closed-fibre components now gives
\[
Q_{Y'}(f^*F_1,\ldots,f^*F_n)=Q_Y(F_1,\ldots,F_n). \tag{PI}
\]
Components whose images have dimension below \(p\) give zero. Extend from effective last slots by differences. This proof never sends local degree through global Chow rational equivalence.

For the raw ideal diagonal let \(D_I=-E_I\). The standard finite Rees algebra has \(\mathcal O_Y(D_I)=\mathcal O_Y(1)\). Complete Stacks 30.14.3 gives \(H^0(Y,\mathcal O_Y(qD_I))=I^q\) for all large \(q\). Complete 30.16.1(4) gives higher cohomology vanishing. The actual Cartier sequence then yields
\[
\ell_R(I^{q-1}/I^q)=\chi_R(E_I,\mathcal O_{E_I}((q-1)D_I)).
\]
Its degree-\(p\) coefficients give \(Q_Y(D_I^{[p]},E_I)=e_R(I)\), or \(-Q_Y(D_I^{[n]})=e_R(I)\), including \(n=1\). No normality, CM property, residue-field infinitude or parameter-colength formula is used.

For \(I_1,\ldots,I_n\), take a common product blowup. The product ideal \(\prod I_j^{t_j}\) has divisor \(\sum t_jD_{I_j}\). Projection (PI) applies its diagonal formula on this common model. If \(H\) is the degree-\(n\) length form, \(e(\prod I_j^{t_j})=n!H(t)\), as seen by substituting exponents \(q t_j\). Equating polynomials on positive integer \(t_j\), then reading \(t_1\cdots t_n\), proves
\[
-Q_Y(D_{I_1},\ldots,D_{I_n})=e_R(I_1,\ldots,I_n). \tag{I}
\]
This is equality of whole polynomials; it makes no coefficientwise positivity inference.
